% GENERATED FILE. Edit canonical lessons, then run npm run generate:books.
% canonical-source-hash: fnv1a64:c60eb9b3021ec0c8
% canonical-lessons: HI-C379-yadi, HI-C379-anyatha, HI-C379-matra, HI-C379-khabardar, HI-C379-patakna

\chapter{Formal Conditions}
\label{ch:hi-formal-conditions}
\hlchaptermodality{379}

\begin{chapteropening}
\textbf{By the end of this chapter:} \emph{I can understand conditions in a formal notice.}

Complete the last lesson of chapter 379: I can understand conditions in a formal notice.
\end{chapteropening}

\section[\texorpdfstring{\hi{यदि} (yadi)}{यदि (yadi)}]{\hi{यदि} (yadi) — if (formal)}

\label{lesson:HI-C379-yadi}



\begin{quote}
Before the new one: say the Hindi for except, then the Hindi for to forbid.
\end{quote}

\subsection*{You'll want to know: \hi{यदि}}
\textbf{\hi{यदि}} — \emph{yadi} — \textquotedblleft{}if (formal)\textquotedblright{}.

\begin{quote}
\hi{यदि} \hi{आप} \hi{सहमत} \hi{हैं} \hi{तो} \hi{कृपया} \hi{यहाँ} \hi{हस्ताक्षर} \hi{कीजिए।}
\end{quote}

\begin{quote}
\emph{yadi āp sahmat haiṁ to kṛpayā yahā̃ hastākṣar kījiye.}
\end{quote}

\begin{quote}
If you agree, please sign here.
\end{quote}

\begin{cousinweb}
\textbf{\hi{यदि}} (\emph{yadi}) is a tatsama, borrowed from Sanskrit \emph{yadi}, \textquotedblleft{}if\textquotedblright{} (Monier-Williams).
\end{cousinweb}

\begin{grammarlens}[title={What you've built}]
One more word for saying what happens if.
\end{grammarlens}

\subsection*{Your turn}
\begin{itemize}
\raggedright
  \item \emph{Say it:} \emph{yadi}
  \item \emph{Say it:} \emph{yadi}, once more
  \item \emph{Say it:} \emph{yadi āp sahmat haiṁ to kṛpayā yahā̃ hastākṣar kījiye.} — If you agree, please sign here.
  \item \emph{Recall:} say the Hindi for except, then the Hindi for to forbid, then say \emph{yadi} again
\end{itemize}

\begin{tcolorbox}[breakable,colback=teal!4,colframe=teal!35!black,title={Before you move on}]
What does \hi{यदि} mean? (\textquotedblleft{}If (formal)\textquotedblright{}.) Say it once more.
\end{tcolorbox}
\section[\texorpdfstring{\hi{अन्यथा} (anyathā)}{अन्यथा (anyathā)}]{\hi{अन्यथा} (anyathā) — otherwise (formal)}

\label{lesson:HI-C379-anyatha}



\begin{quote}
Before the new one: say the Hindi for to forbid, then the Hindi for if (formal).
\end{quote}

\subsection*{You'll want to know: \hi{अन्यथा}}
\textbf{\hi{अन्यथा}} — \emph{anyathā} — \textquotedblleft{}otherwise (formal)\textquotedblright{}.

\begin{quote}
\hi{समय} \hi{पर} \hi{किराया} \hi{दीजिए}, \hi{अन्यथा} \hi{घर} \hi{छोड़ना} \hi{पड़ेगा।}
\end{quote}

\begin{quote}
\emph{samay par kirāyā dījiye, anyathā ghar choṛnā paṛegā.}
\end{quote}

\begin{quote}
Pay the rent on time; otherwise you will have to leave the house.
\end{quote}

\begin{cousinweb}
\textbf{\hi{अन्यथा}} (\emph{anyathā}) is a tatsama, borrowed from Sanskrit \emph{anyathā}, \textquotedblleft{}otherwise, in another way\textquotedblright{} (Monier-Williams).
\end{cousinweb}

\begin{grammarlens}[title={What you've built}]
One more word for saying what happens if.
\end{grammarlens}

\subsection*{Your turn}
\begin{itemize}
\raggedright
  \item \emph{Say it:} \emph{anyathā}
  \item \emph{Say it:} \emph{anyathā}, once more
  \item \emph{Say it:} \emph{samay par kirāyā dījiye, anyathā ghar choṛnā paṛegā.} — Pay the rent on time; otherwise you will have to leave the house.
  \item \emph{Recall:} say the Hindi for to forbid, then the Hindi for if (formal), then say \emph{anyathā} again
\end{itemize}

\begin{tcolorbox}[breakable,colback=teal!4,colframe=teal!35!black,title={Before you move on}]
What does \hi{अन्यथा} mean? (\textquotedblleft{}Otherwise (formal)\textquotedblright{}.) Say it once more.
\end{tcolorbox}
\section[\texorpdfstring{\hi{मात्र} (mātra)}{मात्र (mātra)}]{\hi{मात्र} (mātra) — mere, only (formal)}

\label{lesson:HI-C379-matra}



\begin{quote}
Before the new one: say the Hindi for if (formal), then the Hindi for otherwise (formal).
\end{quote}

\subsection*{You'll want to know: \hi{मात्र}}
\textbf{\hi{मात्र}} — \emph{mātra} — \textquotedblleft{}mere, only (formal)\textquotedblright{}.

\begin{quote}
\hi{यह} \hi{मात्र} \hi{एक} \hi{सवाल} \hi{है}, \hi{अगर} \hi{जवाब} \hi{नहीं} \hi{पता} \hi{तो} \hi{कोई} \hi{बात} \hi{नहीं।}
\end{quote}

\begin{quote}
\emph{yah mātra ek savāl hai, agar javāb nahīṃ patā to koī bāt nahīṃ.}
\end{quote}

\begin{quote}
It's only a question; if you don't know the answer, never mind.
\end{quote}

\begin{cousinweb}
\textbf{\hi{मात्र}} (\emph{mātra}) is a tatsama, borrowed from Sanskrit \emph{mātra}, \textquotedblleft{}a measure; nothing but\textquotedblright{} (Monier-Williams).
\end{cousinweb}

\begin{grammarlens}[title={What you've built}]
One more word for saying what happens if.
\end{grammarlens}

\subsection*{Your turn}
\begin{itemize}
\raggedright
  \item \emph{Say it:} \emph{mātra}
  \item \emph{Say it:} \emph{mātra}, once more
  \item \emph{Say it:} \emph{yah mātra ek savāl hai, agar javāb nahīṃ patā to koī bāt nahīṃ.} — It's only a question; if you don't know the answer, never mind.
  \item \emph{Recall:} say the Hindi for if (formal), then the Hindi for otherwise (formal), then say \emph{mātra} again
\end{itemize}

\begin{tcolorbox}[breakable,colback=teal!4,colframe=teal!35!black,title={Before you move on}]
What does \hi{मात्र} mean? (\textquotedblleft{}Mere, only (formal)\textquotedblright{}.) Say it once more.
\end{tcolorbox}
\section[\texorpdfstring{\hi{ख़बरदार} (khabardār)}{ख़बरदार (khabardār)}]{\hi{ख़बरदार} (khabardār) — beware!, watch out}

\label{lesson:HI-C379-khabardar}



\begin{quote}
Before the new one: say the Hindi for otherwise (formal), then the Hindi for mere.
\end{quote}

\subsection*{You'll want to know: \hi{ख़बरदार}}
\textbf{\hi{ख़बरदार}} — \emph{khabardār} — \textquotedblleft{}beware!, watch out\textquotedblright{}.

\begin{quote}
\hi{ख़बरदार}, \hi{अगर} \hi{तुम} \hi{फिर} \hi{देर} \hi{से} \hi{आए} \hi{तो} \hi{दरवाज़ा} \hi{बंद} \hi{मिलेगा।}
\end{quote}

\begin{quote}
\emph{khabardār, agar tum phir der se āe to darvāzā band milegā.}
\end{quote}

\begin{quote}
Watch out: if you come late again, you'll find the door shut.
\end{quote}

\begin{cousinweb}
\textbf{\hi{ख़बरदार}} (\emph{khabardār}) is usually traced to Persian \emph{khabardār}, \textquotedblleft{}aware, informed\textquotedblright{}. It joins Arabic \emph{khabar}, \textquotedblleft{}news\textquotedblright{}, to Persian \emph{-dār}, \textquotedblleft{}holding\textquotedblright{}: one who has been told.
\end{cousinweb}

\begin{grammarlens}[title={What you've built}]
One more word for saying what happens if.
\end{grammarlens}

\subsection*{Your turn}
\begin{itemize}
\raggedright
  \item \emph{Say it:} \emph{khabardār}
  \item \emph{Say it:} \emph{khabardār}, once more
  \item \emph{Say it:} \emph{khabardār, agar tum phir der se āe to darvāzā band milegā.} — Watch out: if you come late again, you'll find the door shut.
  \item \emph{Recall:} say the Hindi for otherwise (formal), then the Hindi for mere, then say \emph{khabardār} again
\end{itemize}

\begin{tcolorbox}[breakable,colback=teal!4,colframe=teal!35!black,title={Before you move on}]
What does \hi{ख़बरदार} mean? (\textquotedblleft{}Beware!, watch out\textquotedblright{}.) Say it once more.
\end{tcolorbox}
\section[\texorpdfstring{\hi{पटकना} (paṭaknā)}{पटकना (paṭaknā)}]{\hi{पटकना} (paṭaknā) — to slam, to throw down hard}

\label{lesson:HI-C379-patakna}



\begin{quote}
Before the new one: say the Hindi for mere, then the Hindi for beware.
\end{quote}

\subsection*{You'll want to know: \hi{पटकना}}
\textbf{\hi{पटकना}} — \emph{paṭaknā} — \textquotedblleft{}to slam, to throw down hard\textquotedblright{}.

\begin{quote}
\hi{अगर} \hi{तुम} \hi{दरवाज़ा} \hi{पटकोगे} \hi{तो} \hi{बच्चा} \hi{जाग} \hi{जाएगा।}
\end{quote}

\begin{quote}
\emph{agar tum darvāzā paṭkoge to baccā jāg jāegā.}
\end{quote}

\begin{quote}
If you slam the door, the baby will wake up.
\end{quote}

Past, with \hi{ने} (he or she …): \hi{उसने} \hi{दरवाज़ा} \hi{पटका}, \emph{usne darvāzā paṭkā}.

\begin{cousinweb}
\textbf{\hi{पटकना}} (\emph{paṭaknā}) is a tadbhava, inherited from Sanskrit \emph{paṭ}, \textquotedblleft{}a sudden noise\textquotedblright{} (CDIAL 7691). CDIAL glosses the Hindi verb \textquotedblleft{}to dash against\textquotedblright{}.
\end{cousinweb}

\begin{grammarlens}[title={What you've built}]
One more word for saying what happens if.
\end{grammarlens}

\subsection*{Your turn}
\begin{itemize}
\raggedright
  \item \emph{Say it:} \emph{paṭaknā}
  \item \emph{Say it:} \emph{paṭaknā}, once more
  \item \emph{Say it:} \emph{agar tum darvāzā paṭkoge to baccā jāg jāegā.} — If you slam the door, the baby will wake up.
  \item \emph{Recall:} say the Hindi for mere, then the Hindi for beware, then say \emph{paṭaknā} again
\end{itemize}

\begin{tcolorbox}[breakable,colback=teal!4,colframe=teal!35!black,title={Before you move on}]
What does \hi{पटकना} mean? (\textquotedblleft{}To slam, to throw down hard\textquotedblright{}.) Say it once more.
\end{tcolorbox}
